\chapter{System Testing and Verification}

This chapter describes the testing strategy, experimental setup, and benchmark results for the multi-camera tracking system. Testing covers unit and integration validation of the seven-stage pipeline, and end-to-end evaluation against the CityFlowV2 benchmark with comparison to published state-of-the-art methods.

\section{Testing Setup}

\subsection{Hardware}

Model training was performed on Kaggle cloud infrastructure using an NVIDIA Tesla P100 GPU (16~GB HBM2). Pipeline inference and integration testing were conducted on a local development machine with an NVIDIA GTX~1050~Ti (4~GB GDDR5) and an Intel Core~i7 CPU under Windows~11.

\subsection{Dataset}

All experiments use the \textbf{CityFlowV2} dataset from the AI~City Challenge 2022 Track~1. CityFlowV2 contains 46 synchronised cameras deployed at real urban intersections, with 880 globally annotated vehicle identities. The dataset is organised into three scenes:

\begin{table}[htbp]
\centering
\caption{CityFlowV2 dataset composition.}
\label{tab:dataset}
\begin{tabular}{lcc}
\hline
\textbf{Scene} & \textbf{Cameras} & \textbf{Description} \\ \hline
S01 & 3  & Urban intersection cluster \\
S02 & 4  & Urban intersection cluster \\
S03 & 6  & Urban intersection cluster \\
Total & 46 & 880 annotated vehicle IDs \\ \hline
\end{tabular}
\end{table}

Ground-truth annotations use 1-based frame IDs in MOT format. The pipeline internally uses 0-based frame IDs and converts to 1-based only at export time.

\section{Testing Plan and Strategy}

\subsection{Module Testing}

Each pipeline stage has a corresponding pytest test module under \texttt{tests/}. A smoke-test mode processes only 10~frames per camera and validates that all seven stages execute end-to-end without error in under two minutes, confirming output file presence and schema correctness at each stage boundary. Specific validations include:

\begin{itemize}
    \item Stage~0: frame count in manifest matches extracted JPEG count per camera.
    \item Stage~1: each tracklet JSON contains required fields (\texttt{track\_id}, \texttt{camera\_id}, \texttt{class\_name}, \texttt{frames}).
    \item Stage~2: embedding array shape is $N \times 384$; all rows are unit-norm.
    \item Stage~3: FAISS index size equals embedding count; SQLite row count matches.
    \item Stage~4: no global trajectory assigns the same camera's tracklet to two different global IDs (conflict-free constraint).
    \item Stage~5: TrackEval completes and produces HOTA, IDF1, MOTA in output JSON.
\end{itemize}

\subsection{Integration Testing}

Stage-by-stage integration is verified by running consecutive stage pairs on a sample of Scene~S02 (4~cameras, $\approx$2~minutes of footage). Each stage reads its predecessor's disk outputs, confirming that the file format contracts are satisfied across stage boundaries without manual intervention.

\section{Test Schedule}

\begin{table}[htbp]
\centering
\caption{Per-stage milestones.}
\label{tab:schedule}
\begin{tabular}{lll}
\hline
\textbf{Stage} & \textbf{Component} & \textbf{Status} \\ \hline
Stage~0 & Ingestion \& preprocessing & Complete \\
Stage~1 & Detection \& tracking      & Complete \\
Stage~2 & Feature extraction          & Complete \\
Stage~3 & FAISS indexing              & Complete \\
Stage~4 & Cross-camera association    & Complete \\
Stage~5 & Evaluation (TrackEval)      & Complete \\
Stage~6 & Visualisation \& export     & Complete \\ \hline
\end{tabular}
\end{table}

\section{Comparative Results}

\subsection{Multi-Camera Tracking: CityFlowV2}

Table~\ref{tab:mtmc_results} compares our system against published state-of-the-art methods on the CityFlowV2 (AI~City Challenge 2022 Track~1) benchmark. IDF1 is the primary ranking metric.

\begin{table}[htbp]
\centering
\caption{Multi-camera tracking results on CityFlowV2.}
\label{tab:mtmc_results}
\begin{tabular}{llcc}
\hline
\textbf{Method} & \textbf{Dataset} & \textbf{IDF1 (\%)} & \textbf{MOTA (\%)} \\ \hline
AIC22 1st place            & CityFlowV2 & [AIC22-1ST-IDF1]  & [AIC22-1ST-MOTA]  \\
AIC22 2nd (Li et al.)      & CityFlowV2 & 84.37             & ---               \\
AIC21 1st (Liu et al.)     & CityFlowV2 & $\approx$84.1     & ---               \\
Human baseline             & CityFlowV2 & [HUMAN-IDF1]      & [HUMAN-MOTA]      \\
\textbf{Ours (TransReID ViT-B/16 CLIP)} & CityFlowV2 & \textbf{[OUR-IDF1]} & \textbf{[OUR-MOTA]} \\ \hline
\end{tabular}
\end{table}

\begin{figure}[htbp]
\centering
\framebox[0.85\textwidth]{\parbox{0.83\textwidth}{\centering\vspace{2.5cm}\textit{Figure placeholder: IDF1 bar chart --- Ours vs.~AIC22 1st vs.~AIC22 2nd vs.~AIC21 1st vs.~Human}\vspace{2.5cm}}}
% TODO: insert bar chart — IDF1 comparison: Ours vs AIC22 1st vs AIC22 2nd vs AIC21 1st vs Human
\caption{IDF1 comparison of our system against state-of-the-art methods on CityFlowV2.}
\label{fig:mtmc_chart}
\end{figure}

\subsection{Re-Identification Model Performance}

Table~\ref{tab:reid_results} reports Re-ID performance. Literature results for TransReID ViT-Base/16 on standard benchmarks are included for reference.

\begin{table}[htbp]
\centering
\caption{Re-identification results.}
\label{tab:reid_results}
\begin{tabular}{llcc}
\hline
\textbf{Method} & \textbf{Dataset} & \textbf{mAP (\%)} & \textbf{Rank-1 (\%)} \\ \hline
TransReID ViT-B/16 (paper)         & Market-1501  & 89.0  & 95.1  \\
TransReID ViT-B/16 (paper)         & VeRi-776     & 82.1  & 97.4  \\
Human baseline                     & CityFlowV2   & [HUMAN-MAP]   & [HUMAN-R1]   \\
\textbf{Ours (TransReID ViT-B/16 CLIP)} & CityFlowV2 & \textbf{[OUR-MAP]} & \textbf{[OUR-R1]} \\ \hline
\end{tabular}
\end{table}

\begin{figure}[htbp]
\centering
\framebox[0.85\textwidth]{\parbox{0.83\textwidth}{\centering\vspace{2.5cm}\textit{Figure placeholder: TransReID training loss and Rank-1 accuracy curves across epochs on CityFlowV2}\vspace{2.5cm}}}
% TODO: insert figure — TransReID training loss and Rank-1 accuracy curves across epochs on CityFlowV2
\caption{TransReID ViT-Base/16 CLIP training loss and Rank-1 accuracy during fine-tuning on CityFlowV2.}
\label{fig:training_curves}
\end{figure}

\subsection{Association Hyperparameter Ablation}

The association stage (Stage~4) was subjected to an extensive ablation study with over [PLACEHOLDER] configuration variants evaluated on the CityFlowV2 validation set. Key findings:

\begin{itemize}
    \item Similarity threshold: optimal at 0.53 (values below 0.45 merge unrelated tracks; above 0.60 fragment correct identities).
    \item Appearance weight: optimal at 0.70; increasing HSV weight beyond 0.30 degrades IDF1.
    \item Conflict-free constraint: provides consistent improvement over unconstrained connected components.
    \item K-reciprocal re-ranking: no IDF1 improvement on this dataset; disabled in final system.
\end{itemize}

\begin{figure}[htbp]
\centering
\framebox[0.85\textwidth]{\parbox{0.83\textwidth}{\centering\vspace{2.5cm}\textit{Figure placeholder: Ablation study --- IDF1 sensitivity to Stage~4 hyperparameters}\vspace{2.5cm}}}
% TODO: insert bar chart — IDF1 vs key hyperparameter variants (sim_threshold, appearance_weight, algorithm)
\caption{Ablation study: IDF1 sensitivity to key Stage~4 hyperparameters.}
\label{fig:ablation}
\end{figure}

\subsection{Comparison to the AIC22 Winner: Method Completeness and Component Ablation}

The AI~City Challenge 2022 Track~1 winner (CityTrack, Team~28~\cite{yang2022citytrack})
reports 84.86\% IDF1 using a five-model ensemble and seven engineering components.
Our single-model system reaches 77.94\% IDF1, a gap of 6.93~percentage points.
To understand whether this gap is closable through association engineering or is
fundamentally feature-quality-limited, we audited all seven CityTrack components
against our pipeline and re-implemented the three that were genuinely absent.
Table~\ref{tab:citytrack_audit} summarises the audit.

\begin{table}[htbp]
\centering
\caption{Audit of our pipeline against the seven components of the AIC22 Track~1
winner (CityTrack). ``Enabled'' = active in the production configuration;
``Implemented (off)'' = present but disabled because it was found neutral or harmful;
``Added (off)'' = newly implemented for this study behind a default-off flag.}
\label{tab:citytrack_audit}
\begin{tabular}{p{0.32\textwidth}p{0.22\textwidth}p{0.34\textwidth}}
\hline
\textbf{CityTrack component} & \textbf{Status in our system} & \textbf{Notes} \\ \hline
1. Zone-based search-space reduction & Implemented (off) & Soft $\pm0.03$ similarity bonus, not hard pruning; $-0.4$~pp when enabled. \\
2. Spatio-temporal time window & Enabled & Stronger variant: hard gate outside window $+$ Gaussian score inside, vs.\ their $\times 2$ penalty. \\
3. Stationary Sensitive Association (SSA) & Added (off) & Freezes stationary boxes to recent high-confidence detections. \\
4. Trajectory Re-Link (TRL) & Partial (intra on) & Intra-camera ReID-cosine relink enabled; cross-camera motion relink disabled ($-3.8$ to $-13.2$~pp). \\
5. Bidirectional tracking (BT) & Added (off) & Forward $+$ backward passes merged by IoU$+$ReID gating. \\
6. Occlusion-aware distance matrix & Added (off) & Per-tracklet occlusion flag applies a similarity penalty to occluded pairs. \\
7. Box-grained matching $+$ $k$-reciprocal & Implemented (off) & Mean-pooled embeddings $+$ conflict-free connected components by default; $k$-reciprocal and box-grained multi-query both hurt our 280D-PCA features. \\ \hline
\end{tabular}
\end{table}

Four of the seven components were already present (two enabled, two implemented but
disabled after being found harmful with our current features). The three genuinely
missing components --- SSA, bidirectional tracking, and the occlusion-aware distance
matrix --- were re-implemented behind default-off configuration flags and validated
with 84 unit tests, then ablated on a three-camera CityFlowV2 subset (scene S02:
cameras c006/c007/c008). Table~\ref{tab:citytrack_ablation} reports this ablation.
We stress that the S02-subset IDF1 is \emph{not} comparable to the full-dataset
77.94\% headline; it is an internal diagnostic on a hard scene used only to rank the
new components against one another.

\begin{table}[htbp]
\centering
\caption{Component ablation on the CityFlowV2 S02 three-camera subset
(c006/c007/c008). Metric is subset MTMC IDF1; deltas are relative to the
fresh-feature baseline. ID switches (IDsw) are reported as a reliability check.
The standalone occlusion row is shown struck through because it was measured on
\emph{cached} Stage-2 features and is not comparable; the fair occlusion estimate
is (all~$-$~BT) $= -1.03$~pp.}
\label{tab:citytrack_ablation}
\begin{tabular}{lccc}
\hline
\textbf{Configuration} & \textbf{S02 IDF1 (\%)} & \textbf{$\Delta$ (pp)} & \textbf{IDsw} \\ \hline
Baseline (no new components) & 61.79 & ---    & 43  \\
SSA only                     & 61.97 & $+0.18$ & 43  \\
Bidirectional only           & 66.89 & $+5.10$ & 149 \\
SSA $+$ BT $+$ Occlusion     & 65.87 & $+4.08$ & 142 \\
Occlusion only (cached feats) & \textit{68.29} & \textit{($+6.50$)} & \textit{100} \\ \hline
\end{tabular}
\end{table}

The ablation yields three honest verdicts. \textbf{SSA is neutral}: $+0.18$~pp is
within run-to-run noise and the ID-switch count is unchanged (43). \textbf{Bidirectional
tracking is the only component with a sizeable subset gain} ($+5.10$~pp), but it
triples local ID switches (43~$\rightarrow$~149), which means the apparent IDF1 gain
is partly recovered detections offset by new identity fragmentation; it therefore
requires full-dataset confirmation before it can be trusted, and remains default-off.
\textbf{The occlusion-aware matrix is neutral-to-negative under a fair comparison}:
its eye-catching standalone $+6.50$~pp was a measurement artifact --- that run reused
cached Stage-2 features (its ID switches jumped 43~$\rightarrow$~100 even though
occlusion is a pure Stage-4 operation that cannot change upstream tracking), whereas
the baseline recomputed features fresh. On an equal-footing comparison the occlusion
contribution is (all~$-$~BT)~$=65.87-66.89=-1.03$~pp. We deliberately do not claim the
$+6.50$~pp as a result.

Taken together, the audit supports our central finding: the 6.93~pp gap to the AIC22
winner is \emph{not} attributable to missing association machinery. Six of seven
components are present or were added, and none of the three additions produces a
confirmed, feature-fair full-pipeline improvement. The dominant remaining lever is
feature quality --- specifically the gap between our single TransReID ViT-B/16 model
and CityTrack's five-model ensemble --- consistent with the five-axis association
plateau established in the preceding ablation.
